%%%%%%%%%%%%%%%%%%%%%%%%%%%%%%%%%%%%%%%%%%%%%%%%%%%%%%%%%%%%%%%%%%%%%%%%%%%%%%%%
%% 
%% Abstract of your thesis in default document language
%% 

\cleardoubleoddpage%  Make sure to start abstract on a new odd page
\begin{abstract*}
    Scheduling\todo{\emph{personal} scheduling}
    systems require informal and natural language constraints,
    such as preferred times and dependencies, to be formalized in order to conform to a common scheme without ambiguities.
    Current approaches are either difficult to convert into solvable
    constraints\todo{not really; the problem is they don't even let you specify constraints! e.g. Google Calendar}
    or are too complex to be solved efficiently within this scope\todo{what scope?}. \\
    Therefore, we define an intermediate language to bridge the gap between complex high-level
    and solver-ready low-level constraint languages
    If existing languages can define a compilation step into the proposed intermediate language,
    they can be readily translated into constraints that are then supplied to a constraint solver. \\
    The proposed design is independent of any specific solver technology and avoids prescribing how constraints should be solved. Instead, the thesis focuses on defining the structure and semantics of the intermediate language, along with
    guidelines\todo{you do much better than guidelines, you provide \emph{formal semantics}!}
    for validating and comparing solutions, without constraining the methods used to obtain them.
\end{abstract*}

%% 
%%%%%%%%%%%%%%%%%%%%%%%%%%%%%%%%%%%%%%%%%%%%%%%%%%%%%%%%%%%%%%%%%%%%%%%%%%%%%%%%


%%%%%%%%%%%%%%%%%%%%%%%%%%%%%%%%%%%%%%%%%%%%%%%%%%%%%%%%%%%%%%%%%%%%%%%%%%%%%%%%
%% 
%% Abstract of your thesis in an additional language
%% 

% \cleardoubleoddpage%  Make sure to start abstract on a new odd page
% \begin{otherlanguage}{ngerman}
%     \begin{abstract*}
%     \end{abstract*}
% \end{otherlanguage}

%% 
%%%%%%%%%%%%%%%%%%%%%%%%%%%%%%%%%%%%%%%%%%%%%%%%%%%%%%%%%%%%%%%%%%%%%%%%%%%%%%%%